\section{Comparison of Selection Criteria}
\label{sec:selection-ladder}

The main experiments compare DPP with random base selection, but do not
show how much of the gain comes from target conditioning, surrogate
information, or diversity. We therefore compare a broader set of
selection rules while holding the targets, poison budget, downstream
optimizer, and victim runs fixed. We consider uninformed selection
(First-$m$, Random), target-independent difficulty criteria (GraNd,
EL2N, and boundary proximity), increasingly informative
target-conditioned criteria (pixel distance, feature similarity, and
target relevance), and our full pointwise and diversity-aware selectors.
Bottom-$m$, which selects the lowest-scoring examples according to
Eq.~\eqref{eq:selection-score}, is included only as a diagnostic
reference.

All surrogate-based criteria use the same $K=20$ ensemble. GraNd and
EL2N follow \citet{paul2021deep}; Boundary selects by $B_i$, Feature
similarity uses cosine similarity to $x_t$ from a single surrogate,
Target relevance uses $R_i(x_t)$, Greedy selects the top-$m$ examples
according to $B_i+\lambda R_i(x_t)$, and DPP applies the diversity-aware
selection of Section~\ref{sec:methodology}. We evaluate five fixed
bird$\rightarrow$dog targets with five victim runs per target at
$\epsilon=2\times10^{-3}$.

\begin{table}[!ht]
\centering
\caption{
Comparison of base-selection criteria on bird$\rightarrow$dog at
$\epsilon=2\times10^{-3}$. All methods use the same targets, victim
runs, poison budget, and downstream attack settings. Entries report
mean ASR (\%).
}
\label{tab:selection-ladder}
\small
\setlength{\tabcolsep}{8pt}
\renewcommand{\arraystretch}{1.08}
\begin{tabular}{lcccc}
\toprule
&
\multicolumn{2}{c}{\textbf{ConvNetBN}} &
\multicolumn{2}{c}{\textbf{ResNet20BN}} \\
\cmidrule(lr){2-3}
\cmidrule(lr){4-5}
\textbf{Selection} &
\textbf{FC} &
\textbf{GM} &
\textbf{FC} &
\textbf{GM} \\
\midrule
Bottom-$m$        & 0.0 & 0.0 & 0.0 & 12.0 \\
First-$m$         & 0.0 & 0.0 & 8.0 & 24.0 \\
Random            & 0.0 & 4.0 & 4.0 & 16.0 \\
\midrule
GraNd             & 64.0 & 40.0 & 4.0 & 36.0 \\
EL2N              & 60.0 & 52.0 & 12.0 & 36.0 \\
Boundary ($B_i$)  & 52.0 & 48.0 & 12.0 & 40.0 \\
\midrule
Pixel distance    & 44.0 & 40.0 & 4.0 & 36.0 \\
Feature similarity ($K=1$) & 48.0 & 28.0 & 4.0 & 20.0 \\
Target relevance ($R_i$)    & 44.0 & 36.0 & 0.0 & 24.0 \\
\midrule
Greedy            & 60.0 & 36.0 & 4.0 & 36.0 \\
DPP               & 60.0 & 36.0 & 12.0 & 40.0 \\
\bottomrule
\end{tabular}
\end{table}

Table~\ref{tab:selection-ladder} separates the contribution of the
different information sources used for base selection. Bottom-$m$,
First-$m$, and Random quantify variation due to uninformed base choice;
GraNd, EL2N, and Boundary test whether generally difficult or influential
examples are sufficient; and the remaining criteria progressively
incorporate information about the particular target. The comparison
between Target relevance, Greedy, and DPP further isolates the
contribution of boundary information and set-level diversity. Across
the four model--attack settings, Random achieves 6.0\% ASR on average,
compared with 40.0\% for the strongest target-independent criterion
(EL2N), 34.0\% for Greedy, and 37.0\% for DPP.


\paragraph{Effect of perturbation budget.}
We additionally examine whether the importance of base selection changes
with the perturbation radius. Using GM with ConvNetBN, we select each
base set once and reuse it across
$\epsilon\in\{2\times10^{-3},5\times10^{-3},1\times10^{-2}\}$. Thus, only
the admissible perturbation neighborhood changes across columns.

\begin{table}[!ht]
\centering
\caption{
Effect of perturbation budget on base selection using GM with ConvNetBN
on bird$\rightarrow$dog. The selected bases are fixed across budgets.
Entries report mean ASR (\%).
}
\label{tab:selection-ladder-budget}
\small
\setlength{\tabcolsep}{9pt}
\renewcommand{\arraystretch}{1.08}
\begin{tabular}{lccc}
\toprule
\textbf{Selection} &
$\boldsymbol{2{\times}10^{-3}}$ &
$\boldsymbol{5{\times}10^{-3}}$ &
$\boldsymbol{1{\times}10^{-2}}$ \\
\midrule
Bottom-$m$      & 0.0 & 0.0 & 4.0 \\
First-$m$       & 0.0 & 16.0 & 44.0 \\
Random          & 4.0 & 12.0 & 40.0 \\
Pixel distance  & 40.0 & 56.0 & 60.0 \\
DPP             & 36.0 & 64.0 & 60.0 \\
\midrule
DPP $-$ Bottom-$m$ & +36.0 & +64.0 & +56.0 \\
\bottomrule
\end{tabular}
\end{table}

The gap between favorable and unfavorable base choices changes from 36.0
percentage points at $\epsilon=2\times10^{-3}$ to 56.0 at
$\epsilon=1\times10^{-2}$. This directly tests the motivation of
Section~\ref{sec:methodology}: when the perturbation budget is small,
the clean base more strongly restricts the region accessible to the
poison optimizer, whereas this dependence can weaken as the admissible
neighborhood expands.

Overall, these experiments quantify how much attack performance is
determined by base choice before perturbation optimization begins.
They also distinguish the benefit of simply conditioning selection on
the target from the additional gains obtained by the proposed
pointwise score and diversity-aware subset construction.